% titel: Blatt 3, A1: Natriumcarbonat aus NaOH und CO2
% Teil II: der durchgerechnete Weg.

\subsection*{Gegeben und gesucht}
Reaktion: $2\,\mathrm{NaOH(s)} + \mathrm{CO_2(g)} \rightarrow
\mathrm{Na_2CO_3(s)} + \mathrm{H_2O(l)}$. Gegeben
$m(\mathrm{NaOH}) = \qty{100}{g}$, $m(\mathrm{CO_2}) = \qty{60}{g}$. Gesucht
ist die gebildete Masse $m(\mathrm{Na_2CO_3})$.

Molare Massen (gerundet): $M(\mathrm{NaOH}) \approx \qty{40.00}{g/mol}$,
$M(\mathrm{CO_2}) \approx \qty{44.01}{g/mol}$,
$M(\mathrm{Na_2CO_3}) \approx \qty{105.99}{g/mol}$.

\subsection*{Ansatz}
Stoffmengen beider Edukte aus $n=m/M$, dann die Quotienten $n/\nu$ vergleichen
($\nu(\mathrm{NaOH})=2$, $\nu(\mathrm{CO_2})=1$). Das kleinere Verhältnis
gibt das limitierende Reagenz; daraus folgt die Produktmenge.

\subsection*{Rechnung}
\begin{align*}
  n(\mathrm{NaOH}) &= \frac{100}{40.00} = 2.50\ \mathrm{mol}, \qquad
  n(\mathrm{CO_2}) = \frac{60}{44.01} = 1.36\ \mathrm{mol}, \\[2pt]
  \frac{n(\mathrm{NaOH})}{2} &= \frac{2.50}{2} = 1.25, \qquad
  \frac{n(\mathrm{CO_2})}{1} = 1.36.
\end{align*}
Da $1.25 < 1.36$, ist $\mathrm{NaOH}$ das limitierende Reagenz. Aus der
Stöchiometrie folgt pro $2\,\mathrm{NaOH}$
$1\,\mathrm{Na_2CO_3}$, also
\begin{align*}
  n(\mathrm{Na_2CO_3}) &= \frac{1}{2}\,n(\mathrm{NaOH}) = \frac{2.50}{2}
  = 1.25\ \mathrm{mol}, \\
  m(\mathrm{Na_2CO_3}) &= 1.25 \times 105.99
  \approx \qty{132.5}{g} \approx \qty{132}{g}.
\end{align*}

\begin{ergebnis}
Aus $\qty{100}{g}$ $\mathrm{NaOH}$ und $\qty{60}{g}$ $\mathrm{CO_2}$ entstehen
$m(\mathrm{Na_2CO_3}) \approx \qty{132}{g}$;
$\mathrm{NaOH}$ ist der begrenzende Reaktand.
\end{ergebnis}

\begin{kontrolle}
$\mathrm{CO_2}$ ist im Überschuss: verbraucht wird
$1.25\ \mathrm{mol}$ von $1.36\ \mathrm{mol}$, es bleiben $\approx
0.11\ \mathrm{mol}$ \,($\approx \qty{5}{g}$) übrig. Die Masse des Produkts
(132 g) liegt zwischen den beiden Eduktmassen (100 g, 60 g), was plausibel
ist: Aus $1.25\ \mathrm{mol}$ $\mathrm{NaOH}$ entstehen $1.25\ \mathrm{mol}$,
also $\approx 132$ g $\mathrm{Na_2CO_3}$.
\end{kontrolle}

\begin{merksatz}
Limitierendes Reagenz: beide Edukt-Stoffmengen bilden, jeweils durch den
Stöchiometriefaktor teilen; der kleinere Quotient entscheidet und liefert die
mögliche Produktmenge.
\end{merksatz}
